% P2 independent post-outcome exact-group leave-one-out influence results
% from original V29 frozen 49 x 4 x 3 numerical result.
\subsection{Localisation of the observed AUC change}
\label{sec:p2-local-effect}
The complete 49-group average loss is not representative of
the typical held recording. In the source-macro policy,
the affected content group represented by released ID~4
has $Y_{AA}=0.8550$ under the original fit and
$Y_{BB}=0.5925$ under content-disjoint refitting,
a paired difference of $-0.2625$ AUC.
Holding its source model fixed while changing the frozen
selector alone produced $Y_{AB}=0.8400$;
changing the source model while holding the old
$\lambda$ produced $Y_{BA}=0.6075$.
Thus the contrast at this specific held group is
strongly associated with removing same-content
source training records. The cross-fitted model also
changes the tangent representation, so we do not
ascribe the effect solely to memorisation by LDA.

The remaining 48 content groups had a mean paired
$Y_{BB}-Y_{AA}$ difference of $-0.00078125$.
Therefore
\[
\underbrace{-0.00612245}_{\mathrm{all}\ 49}
=
\frac{
\underbrace{-0.262500}_{\mathrm{overlap\ group}}
+48\underbrace{(-0.00078125)}_{\mathrm{other}\ 48}
}{49}.
\]
This identity establishes the \emph{localised}
nature of the study's empirical target contrast.
It is stronger evidence of an evaluation-unit
failure than a claim that purging improves or
reduces AUC by a near-universal constant.

\begin{table}[tb]
\centering\small
\caption{Original 49-held-content-record
paired full-model/full-source-selection contrasts
by policy. No values come from new target recordings.
All effects are AUC differences (not accuracy pp);
the 48-group column removes the content-overlap
target represented by ID~4 from the \emph{average}
only, without refitting the remaining models again.
}
\label{tab:p2-localisation-controls}
\setlength{\tabcolsep}{4.5pt}
\resizebox{\linewidth}{!}{%
\begin{tabular}{p{.37\linewidth}rrrr}
\toprule
Selector & All 49 & ID~4 & Other 48 &
Selected $\lambda$ changed\\
\midrule
Source-macro
& $-0.006122$ & $-0.262500$ & $-0.000781$ & 37/49\\
Proper predictive Gaussian
& $-0.003724$ & $-0.245000$ & $+0.001302$ & 0/49\\
Historical class-centred CovScore
& $-0.006071$ & $-0.260000$ & $-0.000781$ & 0/49\\
\bottomrule
\end{tabular}}
\end{table}
The two covariance selectors never changed their
source-selected $\lambda$ in any of 49 model conditions,
yet the overlapping content group had losses of
$0.2450$ and $0.2600$ AUC when its original
source model was purged. This is a valuable
built-in negative control against an explanation
based entirely on parameter retuning.
Conversely, the source-macro selector changed
its frozen parameter in 37/49 groups, but
this frequent change did not yield a commensurate
aggregate advantage on the other 48 group targets.

\subsection{Person-level heterogeneity and alternative reporting}
\label{sec:p2-heterogeneity}
For the source-macro policy, the 49 matched
full-change values ranged from $-0.2625$
to $+0.0375$ AUC. Their median was exactly
zero, with 24 positive, 21 negative and four
zero changes. Eight content groups lost more
than 0.02 AUC, and six gained more than
0.02 AUC. The overall sign of the average
therefore does not describe all held-group
outcomes. The largest absolute negative
difference occurred in the ID~4 content group;
the next most extreme absolute changes
were substantially smaller (at most
0.0375 in the original frozen row table).
The mean paired change over all possible
48-of-49 leave-one-recording-group descriptive
averages lay between about $-0.00703$
and $-0.00078$ AUC. Removing the
overlap-affected content group makes the
observed mean small, while most other
single-group deletions leave the larger
effect intact.

The nominal 50-ID source-macro mean
$0.594000$ and content-group mean
$0.588673469$ also permit a transparent
reporting-weight check. Because 49
original-model evaluated content representatives
are kept, the 50th repeated released ID
contributed an implied originally reported
AUC of
$50(0.594000)-49(0.588673469)=0.8550$.
This equals the original held outcome for
the overlapping group. The denominator
change alone therefore moves the original
reported mean by $-0.005327$; it must not
be confused with the $-0.006122$ change
on the \emph{matched} 49 original-versus-purged
source models, nor with a patient-level
causal leakage effect. A high-performing
repeat can inflate a macro-average even
without retraining any classifier.

\subsection{Implications for model evaluation workflows}
\label{sec:p2-guidelines}
This case suggests a concrete, low-overhead
procedure for reusing large publicly
released EEG collections:
(1) version-pin the downloaded source files
and the epoch-extraction pipeline,
(2) hash all labelled processed trial contents
under a deterministic byte convention,
(3) assemble connected components of
cross-ID identical contents before any split,
(4) construct training, validation and
held-record partitions on components rather
than released IDs,
(5) fit all preprocessing references,
covariance estimators and hyperparameters
strictly within source content groups,
(6) compute both nominal-ID and unique-content
reporting denominators and
(7) publish any anomalies with non-identifying
reproduction material and an opportunity for
the data steward to clarify provenance.
These steps are safeguards, not a novel
cryptographic algorithm.
A more sophisticated near-duplicate detector
would require a new, separately evaluated
false-positive/false-negative study.

The results also establish an important
negative conclusion: after removing the
overlap group from the descriptive average,
the 48 unaffected held groups exhibited
almost no net source-macro AUC change.
The present evidence supports a local
released-record-content validation risk,
rather than a broad numerical correction
factor transferable to all EEG studies.
